\scanpage{246}
\mspage{192}

signifie donc aussi que l'inclusion des fibres fermées
\[
X_{\unsure{a|0}}\to X,\qquad X_b\to X
\]
sont coasphériques. Mais $X_a\to X$ ne peut être coasphérique que si $X_b=\emptyset$, car pour $x\in X_b\subset X$, $x\backslash X_a$ (ainsi que $X_a/x$) est vide, donc \uline{pas} $W$-asphérique (si on ne suppose $W=\mathrm{Fl}(\mathrm{Cat})$, cas que nous excluons). Donc $\unsure{\Phi|\mathscr{D}}(X)$ n'est propre sur $S$ \uline{que} si $X_a\deleted{\unsure{\text{et}}}\unsure{=}X_b=\emptyset$ (résultat valable pour \uline{\unsure{tte}} $X$ sur $S$, qu'elle soit cofibrée sur $S$ ou $W$-propre, ou non). Si on veut $u,v\in W$, cela implique donc que $X_0=\emptyset$, donc on a prouvé :

\margin{$\vert\!\Vert$}
Soit $X\in\mathrm{Cat}_{/S}$, $S={}$\diagram{\begin{tikzcd}[ampersand replacement=\&,column sep=tiny,row sep=tiny]
\&a\\0\arrow[ur]\arrow[dr]\&\\\&b
\end{tikzcd}}. Si $\Phi(X)$ est propre
sur $S$, alors $X=\emptyset$ (si on suppose $W\ne\mathrm{Fl}(\mathrm{Cat})$).

Mais chemin faisant est apparue une idée bizarre. \unsure{Considérons|Reconsidérons} le \deleted{\unsure{complexe}} diagramme infini, pour $X$ dans $\mathrm{Cat}_{/S}$ \deleted{$\unsure{S}$} :
\[
X\longrightarrow\Phi(X)\longrightarrow\Psi\Phi X\ \unsure{\text{---}|\to}\ \Phi\Psi\Phi(X)\longrightarrow\ldots
\]
Si $X$ est quelconque, ces flèches sont alternativement dans $W_S^d$ et $W_S^g$, donc leurs \unsure{composés|composées} sont dans $W_S$ sans plus. Mais si $X$ \deleted{\unsure{de}} \inserted{est dans} $\underline{\mathrm{Fib}}_W S$, \deleted{\unsure{alors les}} \inserted{comme} cette catégorie est \uline{stable} par $\Phi,\Psi$, et que sur $\underline{\mathrm{Fib}}_W S$ $W_S^d=W_S^g=W_S^f=W_S^{\unsure{\mathrm{univ}|\mathrm{loc}}}$, on voit que dans ce cas les flèches sont toutes dans $W_S^f$. \deleted{En tous cas on} On a envie maintenant de poser $\Phi_\infty(X)=\underrightarrow{\lim}$ de tout ça
% archive sheet number: 245
